%%=============================================================================
%% Conclusie
%%=============================================================================

\chapter{Conclusie}%
\label{ch:conclusie}

Deze bachelorproef onderzocht in welke mate de game-engine Unity geschikt is voor het ontwikkelen van een webgebaseerde simulatie waarin beginnende gebruikers zonder ervaring met wachtwoordmanagers veilig en toegankelijk kunnen oefenen met het gebruik ervan. Op basis van de literatuurstudie, requirementsanalyse, technische evaluatie en het ontwikkelde proof-of-concept kan geconcludeerd worden dat Unity hiervoor een geschikte, maar niet perfecte oplossing is.

\section{Voordelen}
Enerzijds toont dit onderzoek aan dat Unity een duidelijke meerwaarde biedt voor het ontwikkelen van interactieve en visueel ondersteunde leeromgevingen. De gerealiseerde simulatie maakte het mogelijk om gebruikers op een laagdrempelige manier kennis te laten maken met wachtwoordmanagers, zonder gebruik van echte gegevens. De vragenlijst nadien bevestigt dat deze aanpak leidt tot een betere gebruikerservaring, verhoogde bewustwording rond wachtwoordveiligheid en een grotere bereidheid om effectief een wachtwoordmanager te gebruiken. Daarmee levert dit onderzoek een concrete bijdrage aan het domein van cybersecurity-educatie, door aan te tonen dat interactieve simulaties een effectieve aanvulling vormen op traditionele leermethoden.

\section{Beperkingen}
Anderzijds zijn er ook belangrijke beperkingen. Technisch gezien is Unity minder optimaal voor webtoepassingen, wat zich uit in grotere bestanden en langere laadtijden. Daarnaast bleek uit de evaluatie dat niet alle leerdoelen even sterk bereikt werden: concepten zoals hashing en het belang van het masterwachtwoord werden minder goed begrepen door gebruikers. Dit wijst erop dat, hoewel de simulatie effectief is voor algemene bewustwording en praktische vaardigheden, complexe onderliggende beveiligingsconcepten extra ondersteuning en duiding vereisen.

\section{Resultaten}
De bekomen resultaten lagen grotendeels in lijn met de verwachtingen uit de literatuur, waarin interactieve leeromgevingen als veelbelovend worden beschouwd. Toch tonen de resultaten ook aan dat gebruiksvriendelijkheid en begeleiding cruciaal blijven: kleine tekortkomingen in instructies of interface kunnen de leerervaring beïnvloeden. Bovendien was de testgroep beperkt, waardoor de resultaten eerder indicatief zijn en niet meteen veralgemeend kunnen worden.

\section{Toekomstig onderzoek}
Dit onderzoek roept ook nieuwe vragen op die uitnodigen tot verder onderzoek. Zo kan onderzocht worden hoe dergelijke simulaties opgeschaald kunnen worden naar grotere en meer diverse doelgroepen, of hoe ze geïntegreerd kunnen worden in onderwijscontexten. 
Daarnaast is er ruimte om alternatieve technologieën (zoals klassieke webtechnologieën) verder te vergelijken met Unity op vlak van performantie en gebruikservaring.
Verder kan toekomstig onderzoek zich richten op het uitbreiden van de simulatie naar andere cybersecurityconcepten, zoals multifactorauthenticatie (MFA). Aangezien MFA een belangrijke extra beveiligingslaag vormt bovenop wachtwoordgebruik, is het relevant om te onderzoeken hoe gebruikers dit concept begrijpen en toepassen binnen een interactieve leeromgeving. Een simulatie waarin gebruikers bijvoorbeeld werken met een authenticator-app of tijdelijke codes zou kunnen bijdragen aan een beter begrip en een hogere adoptie van deze beveiligingsmaatregel. 
Tot slot kan toekomstig onderzoek zich richten op het verbeteren van de didactische aanpak binnen de simulatie, bijvoorbeeld door complexere concepten beter visueel of interactief uit te leggen.

\section{Besluit}
Samenvattend toont deze bachelorproef aan dat een webgebaseerde simulatie, ontwikkeld in Unity, een waardevolle en haalbare oplossing is voor het ondersteunen van cybersecurity-educatie bij beginnende gebruikers. Hoewel er nog optimalisaties mogelijk zijn, vormt deze aanpak een veelbelovende stap richting effectievere en praktijkgerichte leeromgevingen voor veilig digitaal gedrag.

% TODO: Trek een duidelijke conclusie, in de vorm van een antwoord op de 
% onderzoeksvra(a)g(en). Wat was jouw bijdrage aan het onderzoeksdomein en
% hoe biedt dit meerwaarde aan het vakgebied/doelgroep? 
% Reflecteer kritisch over het resultaat. In Engelse teksten wordt deze sectie 
% `Discussion'' genoemd. Had je deze uitkomst verwacht? Zijn er zaken die nog 
% niet duidelijk zijn? 
% Heeft het onderzoek geleid tot nieuwe vragen die uitnodigen tot verder onderzoek?